\section{Dataset}
\label{sec:dataset}

\subsection{Sourcing}

The dataset used in this study is the \textbf{STRAMPN Histopathological Images for Ovarian Cancer Prediction}, published on IEEE DataPort and accessible at:

\begin{center}
\url{https://ieee-dataport.org/documents/strampn-histopathological-images-ovarian-cancer-prediction}
\end{center}

The dataset consists of histopathological JPEG images drawn from ovarian tissue specimens and is labelled for binary classification: \emph{Ovarian Cancer} versus \emph{Ovarian Non-Cancer}. The dataset was accessed in May 2026. In accordance with IEEE DataPort terms, the dataset is not redistributed as part of this repository; all rights remain with the original authors.

\subsection{Composition}

The full dataset contains 987 images distributed across two classes as shown in Table~\ref{tab:class-composition}.

\begin{table}[H]
  \centering
  \caption{Class composition of the STRAMPN dataset.}
  \label{tab:class-composition}
  \begin{tabular}{lrr}
    \toprule
    \rowcolor{tableheader}
    \textbf{Class} & \textbf{Count} & \textbf{Percentage} \\
    \midrule
    Ovarian Cancer     & 481 & 48.73\% \\
    Ovarian Non-Cancer & 506 & 51.27\% \\
    \midrule
    \textbf{Total}     & \textbf{987} & \textbf{100\%} \\
    \bottomrule
  \end{tabular}
\end{table}

The class distribution is near-balanced, with a modest majority of non-cancer images (51.27\%). This near-parity reduces the risk of class-imbalance bias and justifies the use of accuracy as a supplementary metric alongside AUC-ROC.

\subsection{Train / Validation / Test Split}

A two-stage stratified split was applied to partition the dataset. In the first stage, 75\% of images per class were assigned to a combined train-and-validation pool, with the remaining 25\% reserved as the held-out test set. In the second stage, the train-and-validation pool was further divided 70\%/30\%, yielding approximately 52.5\% of the full dataset for training and 22.5\% for validation.

Both split stages used \texttt{sklearn.model\_selection.train\_test\_split} with \texttt{stratify=labels} and \texttt{random\_state=42}, ensuring class proportions are preserved in every subset and the split is fully reproducible. The resulting partition is shown in Table~\ref{tab:split}.

\begin{table}[H]
  \centering
  \caption{Stratified train / validation / test split.}
  \label{tab:split}
  \begin{tabular}{lrrrr}
    \toprule
    \rowcolor{tableheader}
    \textbf{Split} & \textbf{Cancer} & \textbf{Non-Cancer} & \textbf{Total} & \textbf{\% of Dataset} \\
    \midrule
    Train      & 253 & 265 & 518 & 52.48\% \\
    Validation & 108 & 114 & 222 & 22.49\% \\
    Test       & 120 & 127 & 247 & 25.03\% \\
    \midrule
    \textbf{Total} & \textbf{481} & \textbf{506} & \textbf{987} & \textbf{100\%} \\
    \bottomrule
  \end{tabular}
\end{table}

The fixed \texttt{random\_state=42} is applied consistently throughout \texttt{src/preprocessing.py} and is never altered between the 84 training runs, guaranteeing that every configuration trains on identical data partitions.

\subsection{Data Integrity Verification}

Following the file-copy step, \texttt{src/preprocessing.py} performs a cross-split leakage check by comparing image filename stems across all three splits. Any filename appearing in more than one split is flagged as a violation. The preprocessing output confirmed:

\begin{lstlisting}[style=python, basicstyle=\scriptsize\ttfamily]
PREPROCESSING COMPLETE -- No integrity violations found
\end{lstlisting}

No images are shared between the training, validation, and test directories, confirming that evaluation results are free from data leakage.
